% ---------------------------------------------------------------------
%  CHƯƠNG 2 — CƠ SỞ LÝ THUYẾT VÀ CÔNG NGHỆ
% ---------------------------------------------------------------------
\chapter{Cơ sở lý thuyết và công nghệ}
\label{ch:coso}

\begin{muctieu}
Chương này trình bày các kiến thức nền tảng gồm bài toán dự báo chuỗi thời
gian, các phương pháp trí tuệ nhân tạo được dùng, cách đo sai số, cùng danh
mục công cụ và thư viện để xây dựng hệ thống.
\end{muctieu}

\section{Kiến thức nền tảng}

\subsection{Dự báo quỹ đạo bão như bài toán chuỗi thời gian}

\begin{dinhnghia}[Dự báo quỹ đạo nhiều mốc]
Cho chuỗi quan trắc $\mathbf{x}_{t-T+1}, \dots, \mathbf{x}_t$ của một cơn bão,
mỗi $\mathbf{x}_i \in \mathbb{R}^F$ là vec-tơ đặc trưng tại thời điểm $i$.
Bài toán là học hàm $f$ ánh xạ chuỗi này thành vị trí tâm bão
$\hat{\mathbf{p}}_{t+h} = (\hat\varphi_{t+h}, \hat\lambda_{t+h})$ tại các mốc
$h \in \mathcal{H} = \{6, 12, 24, 48, 72\}$ giờ, với $\varphi$ là vĩ độ và
$\lambda$ là kinh độ.
\end{dinhnghia}

Cơn bão không ``chọn'' đường đi như tác tử trong bài toán tìm đường; nó bị
dòng khí quyển quy mô lớn (dòng dẫn đường, \en{steering flow}) đẩy đi. Vì vậy
``gợi ý đường đi'' ở đây là bài toán \term{hồi quy nhiều đầu ra}: ước lượng vị
trí tương lai và độ bất định của nó.

\subsection{Sai số khoảng cách vòng lớn}

Sai số dự báo được đo bằng khoảng cách trên mặt cầu Trái Đất (công thức
Haversine) giữa vị trí dự báo và vị trí thực:
\begin{equation}
d = 2R \arcsin\sqrt{\sin^2\frac{\Delta\varphi}{2}
    + \cos\varphi_1\cos\varphi_2\sin^2\frac{\Delta\lambda}{2}},
\label{eq:haversine}
\end{equation}
với $R = 6371$~km là bán kính Trái Đất. Đây là độ đo chuẩn trong đánh giá dự
báo quỹ đạo bão, có đơn vị km nên dễ hiểu với người dùng.

\subsection{Các phương pháp mốc so sánh}

\term{Persistence} giả định bão giữ nguyên vận tốc của 6 giờ gần nhất.
\term{CLIPER} \en{(CLImatology and PERsistence)}~\cite{Neumann1972} kết hợp quán
tính với khí hậu học; đồ án dùng một phiên bản rút gọn (CLIPER-proxy) với
persistence có suy giảm. Một phương pháp AI chỉ có giá trị nếu thắng được hai
mốc này.

\subsection{Lập luận dựa trên tình huống: k láng giềng gần nhất}

Phương pháp \term{analog} là dạng lập luận dựa trên tình huống
\en{(case-based reasoning)}~\cite{RussellNorvig2021}: tìm trong thư viện các
tình huống bão lịch sử giống tình huống hiện tại nhất rồi dùng diễn biến tiếp
theo của chúng để dự báo. Khoảng cách giữa hai tình huống $q$ và $c$ là khoảng
cách Euclid trên các nhóm đặc trưng $g$ (đoạn quỹ đạo, vị trí, mùa) đã chuẩn
hóa, mỗi nhóm có trọng số $w_g$ và chia cho căn bậc hai số chiều $n_g$ của nhóm
để các nhóm so sánh được với nhau:
\begin{equation}
D(q, c) = \sqrt{\sum_{g} \frac{w_g^2}{n_g} \,\bigl\lVert \mathbf{z}_g(q) - \mathbf{z}_g(c) \bigr\rVert_2^2} .
\label{eq:knn-dist}
\end{equation}
Dự báo là trung bình có trọng số nghịch đảo khoảng cách của $k$ láng giềng:
$\hat{\mathbf{y}} = \sum_{j=1}^{k} \alpha_j \mathbf{y}_j$, với
$\alpha_j \propto 1 / (D_j + \epsilon)$. Ưu điểm lớn của k-NN là giải thích
được: hệ thống nêu tên các cơn bão lịch sử đã dùng.

\subsection{Gradient boosting}

Gradient boosting~\cite{Friedman2001} xây dựng tuần tự nhiều cây quyết định,
mỗi cây học phần dư của tổ hợp các cây trước. Đồ án dùng bộ hồi quy
\en{histogram gradient boosting} của scikit-learn (300 vòng, tốc độ học 0,05,
tối đa 31 lá mỗi cây), mỗi mốc và mỗi tọa độ một mô hình, đầu vào là cửa sổ
đặc trưng đã trải phẳng.

\subsection{Mạng LSTM}

LSTM \en{(Long Short-Term Memory)}~\cite{Hochreiter1997} là mạng nơ-ron hồi quy
có các cổng quên $f_t$, cổng vào $i_t$ và cổng ra $o_t$ điều khiển ô nhớ $c_t$:
\begin{equation}
\begin{aligned}
c_t &= f_t \odot c_{t-1} + i_t \odot \tanh(W_c[\mathbf{h}_{t-1}, \mathbf{x}_t] + b_c),\\
\mathbf{h}_t &= o_t \odot \tanh(c_t),
\end{aligned}
\end{equation}
nhờ đó giữ được thông tin qua nhiều bước. Trạng thái ẩn cuối cùng
$\mathbf{h}_t$ được đưa qua tầng tuyến tính để sinh ra dự báo cho cả năm mốc
cùng lúc. Biến thể \term{Attention-LSTM}~\cite{Bahdanau2015} thay $\mathbf{h}_t$
bằng tổng có trọng số của mọi trạng thái ẩn, trọng số do cơ chế attention học.

\subsection{Transformer}

Transformer~\cite{Vaswani2017} bỏ cơ chế hồi quy, dùng self-attention để mỗi
bước thời gian ``nhìn'' trực tiếp mọi bước khác:
\begin{equation}
\mathrm{Attention}(Q, K, V) = \mathrm{softmax}\!\left(\frac{QK^\top}{\sqrt{d_k}}\right) V,
\end{equation}
cộng thêm mã hóa vị trí để giữ thứ tự thời gian. Đồ án dùng phần encoder.

\subsection{Huấn luyện và kiểm định}

Mạng nơ-ron được tối ưu bằng Adam~\cite{Kingma2015} với hàm mất mát Smooth~L1,
ít nhạy với ngoại lai hơn bình phương sai số. Để so sánh hai phương pháp,
đồ án dùng \term{bootstrap}~\cite{Efron1993} lấy mẫu lại theo cơn bão, vì các
dự báo của cùng một cơn bão tương quan với nhau; khoảng tin cậy 95\% không
chứa 0 nghĩa là khác biệt có ý nghĩa thống kê.

\begin{ghinho}
Với bài toán dự báo, phải chia tập theo thời gian và giữ nguyên cả cơn bão.
Nếu chia ngẫu nhiên theo điểm quan trắc, mô hình sẽ ``thấy'' các điểm khác của
cùng cơn bão khi huấn luyện, và sai số báo cáo sẽ thấp hơn thực tế.
\end{ghinho}

\section{Công cụ và thư viện}

\begin{table}[H]
\centering
\caption{Công cụ và thư viện sử dụng trong đồ án.}
\label{tab:congcu}
\begin{tabularx}{\textwidth}{>{\bfseries\color{daunavy}}l >{\color{daumaroon}}l X}
\toprule
\rowcolor{gray100}\bfseries\color{daunavy}Thành phần & \bfseries\color{daumaroon}Phiên bản & \leavevmode\bfseries\color{daunavy}Vai trò \\
\midrule
Python        & 3.11 & Ngôn ngữ lập trình chính \\
PyTorch       & 2.x  & Huấn luyện LSTM, Attention-LSTM, Transformer \\
scikit-learn  & 1.x  & Gradient boosting, k-NN, chuẩn hóa \\
pandas, NumPy & 2.x  & Xử lý dữ liệu quỹ đạo dạng bảng \\
xarray        & ---  & Đọc dữ liệu tái phân tích NCEP/NCAR \\
FastAPI       & ---  & API phục vụ dự báo \\
Leaflet       & ---  & Bản đồ web hiển thị quỹ đạo và nón \\
pytest        & ---  & Kiểm thử tự động \\
\bottomrule
\end{tabularx}
\end{table}

\section{Các công trình liên quan}

Dự báo quỹ đạo bão truyền thống đi theo ba hướng. Hướng thống kê với CLIPER
\cite{Neumann1972} là mốc so sánh chuẩn suốt nhiều thập kỷ. Hướng mô hình số
trị khí quyển cho độ chính xác cao nhất nhưng cần siêu máy tính. Hướng học máy
gần đây dùng mạng hồi quy và Transformer học trực tiếp từ dữ liệu best-track
như IBTrACS~\cite{Knapp2010}, đôi khi bổ sung trường khí quyển từ tái phân tích
\cite{Kalnay1996}. Đồ án kế thừa hướng học máy, nhưng khác ở chỗ so sánh công
bằng nhiều họ phương pháp trên cùng tập kiểm thử có kiểm định thống kê, tập
trung vào nhóm bão ảnh hưởng Việt Nam, và dùng k-NN analog để giải thích gợi ý.
